% QUESTAO
Sejam duas matrizes $A$ e $B$, de mesma ordem. Ao adicionarmos os elementos correspondentes (de mesma posição) das matrizes $A$ e $B$, estamos adicionando essas matrizes, ou seja, calculando $A+B$. Assim, dada as matriz $A = \left[ \begin{array}{cc}
1 & 0 \\
-2 & 4
\end{array} \right]$ e $B = \left[ \begin{array}{cc}
8 & -2 \\
-5 & 1
\end{array} \right]$, determine a soma $A^t+B$.

% ALTERNATIVAS
$A^t+B = \left[ \begin{array}{cc} 9 & -4 \\ -5 & 5 \end{array} \right]$
$A^t+B = \left[ \begin{array}{cc} 8 & -4 \\ -5 & 1 \end{array} \right]$
$A^t+B = \left[ \begin{array}{cc} 9 & 0  \\ -5 & 5 \end{array} \right]$
$A^t+B = \left[ \begin{array}{cc}-8 & -4 \\ -5 & 5 \end{array} \right]$
$A^t+B = \left[ \begin{array}{cc} 0 & -5 \\ -4 & 3 \end{array} \right]$


